%% ma: L12-B15-0003
%% lop: 12
%% chuong: 5
%% bai: 15
%% dang: 12.15.2
%% loai: TN4
%% muc: NB
%% dap_an: "A"
%% nguon: "(NBV)-ĐỀ SỐ 3-MỖI NGÀY 1 ĐỀ THI 2025.pdf (D03, MỖI NGÀY 1 ĐỀ - Đề số 3), Phần I Câu 12"
%% kien_thuc: "Bài 15 (phương trình tham số của đường thẳng, vectơ chỉ phương)"
%% trang_thai: chinh-thuc
%% ghi_chu: "Đáp án gốc A đúng. Gần ý tưởng với 12.15.2 (đọc vectơ chỉ phương từ phương trình chính tắc); tách riêng vì đề cho dạng tham số, có thể gộp khi duyệt. Giáo viên duyệt gộp dạng 12.15.3 vào 12.15.2 (10/10/2026)"
\begin{ex}
Trong không gian $Oxyz$, đường thẳng $d$ có phương trình tham số
\[\begin{cases}x=-1+2t\\y=3-t\\z=2+t\end{cases}\]
có một vectơ chỉ phương là
\choice{\True $\vec u=(2;-1;1)$}{$\vec v=(-1;3;2)$}{$\vec a=(-1;2;3)$}{$\vec b=(-1;-1;1)$}
\loigiai{Đường thẳng có phương trình tham số $x=x_0+at$, $y=y_0+bt$, $z=z_0+ct$ nhận $(a;b;c)=(2;-1;1)$ làm một vectơ chỉ phương.}
\end{ex}
